\documentclass[10pt,a4paper]{amsart}
\usepackage[margin=19mm]{geometry}
\usepackage{amsmath,amssymb,mathtools}
\usepackage[scr=rsfs,cal=euler]{mathalfa}
\usepackage{xfrac}
\usepackage[unicode,hidelinks,bookmarks=false]{hyperref}
\newcommand{\ie}{i.e.\ }
\newcommand{\cf}{cf.\ }
\newcommand{\resp}{respectively\ }
\newcommand{\Subsection}{\subsection}
\providecommand{\nobreakdash}{\nobreak\hbox{-}}
\setlength{\emergencystretch}{2em}
\multlinegap0pt
\usepackage{bm}
\usepackage{bbm}
\usepackage{dsfont}
\usepackage{xspace}
\usepackage{cancel}
\usepackage{mathtools}
\usepackage{amsthm}
\makeatletter
\@ifundefined{theorem}{\newtheorem{theorem}{Theorem}}{}
\@ifundefined{lemma}{\newtheorem{lemma}[theorem]{Lemma}}{}
\makeatother
\title[Primes in Arithmetic Progressions to Large Moduli and Siegel]{Primes in Arithmetic Progressions to Large Moduli and Siegel Zeroes}
\author{Thomas Wright}
\date{Source: arXiv:2507.10780v4 (CC BY 4.0)}
\begin{document}
\maketitle

\noindent\emph{Source and licence.} Primes in Arithmetic Progressions to Large Moduli and Siegel Zeroes. Version arXiv:2507.10780v4 (CC BY 4.0), distributed under the Creative Commons Attribution 4.0 International licence, as recorded in the arXiv licence metadata for this exact version and shown on the abstract page https://arxiv.org/abs/2507.10780v4. Full rights notice: licenses/arxiv-2507.10780-CC-BY-4.0.txt.\par\smallskip

\noindent\textit{Editorial scope.} Counted unit: the Theorem whose source label is \texttt{None} in the article (source0.tex lines 630--634, source label \texttt{None}), delivered together with the article's own complete original proof of it (source0.tex lines 636--684, one \texttt{proof} environment of 49 lines). No mathematics is written, repaired or re-derived: statement and proof are the article's own text line for line. Required context retained uncounted: Mun (lines 234--240); lambdaap (lines 266--269); DW (lines 442--450); SWx (lines 500--502); DWLambda (lines 507--509). Every definition, display and label of those lines is retained unchanged. Omitted from this package and not redistributed: 1--233, 241--265, 270--441, 451--499, 503--506, 510--629, 635--635, 685--728. Nothing inside a retained excerpt is omitted or altered: every retained body is delivered exactly as the article's lines are written, so no omission receipt of any kind accompanies this package. Citations: the retained excerpts carry no citation command at all, so no bibliography entry of the article is transcribed and no reference is redistributed. Reference adaptation: none was needed and none was made; every reference inside the delivered unit resolves inside it, so no unresolved reference or citation is produced. Printed slips kept literal: no printed word, symbol, hypothesis or number of the counted statement or of its proof was altered. Layout and class: the article's own class is \texttt{amsart} with packages including latexsym, amsmath, amsfonts, epsfig, graphics, graphicx, amsthm, amssymb; the wrapper is the lane's pinned \texttt{amsart} editorial wrapper at 10pt on A4, which restates the theorem-like environments the retained text needs and adds the article's own macro definitions that the retained text needs (none), with extra wrapper packages (bm, bbm, dsfont, xspace, cancel, mathtools). The delivered numbering is the unit's own. Assets: no retained span contains a figure environment, a graphics-inclusion command, a PDF-page inclusion, a table or a TikZ picture; no image file is copied into this package, the package declares no asset dependency and no image file is redistributed with it. Licence: version arXiv:2507.10780v4 of this article is distributed under the Creative Commons Attribution 4.0 International licence, as recorded in the arXiv licence metadata for this exact version and shown on the abstract page https://arxiv.org/abs/2507.10780v4. A full rights notice is delivered at licenses/arxiv-2507.10780-CC-BY-4.0.txt. Characters: every retained excerpt is pure ASCII, so the delivered engine, XeLaTeX with its default Latin Modern text font, reports no missing character.
\begin{lemma}[Munshi, 2011]\label{Mun}
For any natural number $n$ and any real number $r>2$,
$$\tau(n)\ll \sum_{\substack{d|n \\ d\leq n^\frac 1r}}\tau(d)^\beta,$$
where
$$\beta=-\frac{\log r}{\log 2}+r\left(1+\left(1-\frac 1r\right)\frac{\log\left(1-\frac 1r\right)}{\log 2}\right).$$
In particular, if $r\leq 4$ then $\beta<1$.
\end{lemma}

\begin{lemma}\label{lambdaap} For any $\alpha>0$,
$$S(x,q,a)=\frac{1+\chi_D\left(\frac{aD}{(D,q)}\right)}{\phi(q)}S(x,q)+O\left((Dq)^{\frac 12+\alpha}\right).$$
This is non-trivial if $Dq<x^{\frac 23-3\alpha}$.
\end{lemma}

\begin{lemma}\label{DW}
Let $\sqrt x\leq q<D^{-1}x^{\frac 23-3\alpha}$.  Then for any $\alpha>0$,
$$S_W'(x,q,a)=S_R'(x,q,a)+O\left(q^{\frac 12+\alpha}+\frac{x\log x}{qR^{1-2\alpha}}\right),$$
and
$$S_W'(x,q)=S_R'(x,q)+O\left(\frac{x}{R^{1-2\alpha}}\right).$$
Hence,
$$S_W'(x,q,a)=\frac{1-\chi_D\left(\frac{aD}{(D,q)}\right)}{\phi(q)}S_W'(x,q)+O\left((Dq)^{\frac 12+\alpha}+\frac xqe^{-\frac{\varepsilon\log x}{24\log R}}+\frac{x\log^5 xL(1,\chi)}{\phi(q)}\right).$$

\end{lemma}

\begin{lemma}\label{SWx}
$$S_W'(x,q)=x+O\left(R\log^2 x+x\log xL(1,\chi)\right).$$
\end{lemma}

\begin{gather}\label{DWLambda}
S_W'(x,q)=\sum_{R<n_1\leq x}\lambda'_W\left(n_1\right)+\sum_{\substack{R<p_{-1}\leq x }}\log\left(p_{-1}\right)+\sum_{\substack{n_1n_{-1}\leq x \\ n_1,n_{-1}>R }}\lambda_W\left(n_1\right) \lambda_W'\left(n_{-1}\right).
\end{gather}

\begin{theorem}
Let $\sqrt x<Q<D^{-1}x^{\frac 23-\alpha}$ for any $\alpha>0$.  Then for any fixed integer $a\neq 0$,
$$\sum_{\substack{q\sim Q \\ (a,q)=1}}\left|\psi(x,q,a)-\frac{\psi(x)}{\phi(q)}\right|\ll xe^{-\frac{\alpha\log x}{24\log R}
}+x\log^5 xL(1,\chi)+R\log^2 x.$$
\end{theorem}

\begin{proof}
We again recall the decomposition in (\ref{DWLambda}):
\begin{gather}
S_W'(x,q,a)=\sum_{\substack{R<n_1\leq x \\ n_1\equiv a\pmod q}}\lambda'_W\left(n_1\right)+\sum_{\substack{R<p_{-1}\leq x \\ p_{-1}\equiv a\pmod q}}\log\left(p_{-1}\right)+\mathop{\sum\sum}\limits_{\substack{n_1n_{-1}\leq x \\ n_1,n_{-1}>R \\ n_1n_{-1}\equiv a\pmod q}}\lambda_W\left(n_1\right) \lambda_W'\left(n_{-1}\right).
\end{gather}
We show first that when summed over $q\sim Q$, $S_W'(x,q,a)$ will generally be a good approximation for $\psi(x,q,a)$.  To this end,
\begin{align*}
\sum_{\substack{q\sim Q  \\ (a,q)=1}}&\left|S_W'(x,q,a)-\psi(x,q,a)\right|\\
\leq &\sum_{\substack{q\sim Q  \\ (a,q)=1}}\left(\sum_{\substack{R<n_1\leq x \\ n_1\equiv a\pmod q}}\left(\lambda'_W\left(n_1\right)-\Lambda(n_1)\right)+\mathop{\sum\sum}\limits_{\substack{n_1n_{-1}\leq x \\ n_1,n_{-1}>R \\ n_1n_{-1}\equiv a\pmod q}}\lambda_W\left(n_1\right) \lambda_W'\left(n_{-1}\right)+O\left(x^\frac 14\right)\right)
\end{align*}
Recall that $\lambda'_W\left(n_1\right)-\Lambda(n_1)\geq \lambda(n_1)\log x$.  So by Lemma \ref{lambdaap},
\begin{align*}
\sum_{\substack{q\sim Q  \\ (a,q)=1}}\sum_{\substack{R<n_1\leq x \\ n_1\equiv a\pmod q}}\left(\lambda'_W\left(n_1\right)-\Lambda(n_1)\right)\ll &\log x\sum_{\substack{q\sim Q  \\ (a,q)=1}}\sum_{\substack{R<n_1\leq x \\ n_1\equiv a\pmod q}}\lambda\left(n_1\right)\\
\ll &\log x\sum_{\substack{q\sim Q  \\ (a,q)=1}}\frac{1}{\phi(q)}\sum_{\substack{R<n_1\leq x }}\lambda\left(n_1\right)\\
\ll &xL(1,\chi)\log x.
\end{align*}
Hence
\begin{align*}
\sum_{\substack{q\sim Q  \\ (a,q)=1}}&\left|S_W'(x,q,a)-\psi(x,q,a)\right|\ll &\log x\sum_{\substack{q\sim Q  \\ (a,q)=1}}\left(\mathop{\sum\sum}\limits_{\substack{n_1n_{-1}\leq x \\ n_1,n_{-1}>R \\ n_1n_{-1}\equiv a\pmod q}}\lambda_W\left(n_1\right) \right)+O\left(xL(1,\chi)\log x+Qx^\frac 14\right).
\end{align*}
Now, if $n_1n_{-1}\equiv a \pmod q$ then $q|n_1n_{-1}-a$, and hence we can write
\begin{align*}
\sum_{\substack{q\sim Q  \\ (a,q)=1}}\mathop{\sum\sum}\limits_{\substack{n_1n_{-1}\leq x \\ n_1,n_{-1}>R \\ n_1n_{-1}\equiv a\pmod q}}\lambda_W\left(n_1\right) \leq \mathop{\sum\sum}\limits_{\substack{n_1n_{-1}\leq x \\ n_1,n_{-1}>R }}\lambda_W\left(n_1\right) \tau(n_1n_{-1}-a).
\end{align*}
By Lemma \ref{Mun}, we can bound this with
\begin{align*}
\ll &\sum_{d\leq x^{\frac 13-\alpha}} \tau(d)\mathop{\sum\sum}\limits_{\substack{n_1n_{-1}\leq x \\ n_1,n_{-1}>R \\ n_1n_{-1}\equiv a\pmod d}}\lambda_W\left(n_1\right) .
\end{align*}
Note that if $n_{-1}\geq x^\frac 14>d^{1+\alpha}$ then
$$\sum_{\substack{x^{\frac 13}\leq n_{-1}\leq \frac x{n_1} \\ n_{-1}\equiv a\overline{n_1}\pmod d}}1\ll \frac 1d\sum_{\substack{n_{-1}\leq \frac x{n_1} }}1,$$
while if $n_{-1}<x^\frac 14$ then $n_{1}\geq x^\frac 23$, and hence
$$\sum_{\substack{x^{\frac 23}\leq n_1\leq \frac x{n_{-1}} \\ n_{1}\equiv a\overline{n_{-1}}\pmod d}}\lambda_W(n_1)\ll \sum_{\substack{x^{\frac 23}\leq n_1\leq \frac x{n_{-1}} \\ n_{1}\equiv a\overline{n_{-1}}\pmod d}}\lambda(n_1)\ll  \frac 1{\phi(d)}\sum_{\substack{n_1\leq \frac x{n_{-1}} }}\lambda(n_1).$$
So
\begin{align*}
\sum_{d\leq x^{\frac 13-\alpha}} &\tau(d)\mathop{\sum\sum}\limits_{\substack{n_1n_{-1}\leq x \\ n_1,n_{-1}>R \\ n_1n_{-1}\equiv a\pmod d}}\lambda_W\left(n_1\right) \\
\ll &\sum_{d\leq x^{\frac 13-\alpha}} \frac{\tau(d)}{\phi(d)}\mathop{\sum\sum}\limits_{\substack{n_1n_{-1}\leq x \\ n_1,n_{-1}>R }}\lambda\left(n_1\right) \\
\ll &x\sum_{d\leq x^{\frac 13-\alpha}} \frac{\tau(d)}{d}\left[L(1,\chi)\mathop{\sum}\limits_{\substack{ R<n_{-1}\leq 2x^\frac 23}}\frac{1}{n_{-1}}+\mathop{\sum}\limits_{\substack{R<n_{-1}\leq 2x^\frac 14}}\frac{\lambda(n_1)}{n_{1}}\right]\\
\ll &xL(1,\chi)\log^{3} x.
\end{align*}
Hence
\begin{align*}
\sum_{\substack{q\sim Q  \\ (a,q)=1}}&\left|S_W'(x,q,a)-\psi(x,q,a)\right|\ll xL(1,\chi)\log^4 x+Qx^\frac 14\ll xL(1,\chi)\log^4 x.
\end{align*}
From Lemma \ref{DW}, if $Q<D^{-1}x^{\frac 23-3\alpha}$ then
$$\sum_{\substack{q\sim Q  \\ (a,q)=1}}\left|S_W'(x,q,a)-\frac{1}{\phi(q)}S_W'(x,q)\right|\ll xe^{-\frac{\alpha\log x}{24\log R}}+x\log^5 xL(1,\chi),$$
and we recall from Lemma \ref{SWx} that
$$S'_W(x,q)=\psi(x)+O\left(R\log^2 x+x\log xL(1,\chi)\right).$$
Choosing the largest terms from these expressions, the theorem then follows.
\end{proof}

\end{document}
